% Frames cut from session1_slides.tex on 2026-10-26 (old Session 1 Part B exercise).
% Destination: session2_slides.tex, block 17 of admin/course_plan_4_sessions.md
% ("Exercise: build the county panel in plan mode, evaluate the plan, debrief", 25 min).
% Not compiled on its own. Insert these frames into the Session 2 deck, then delete
% this file.
%
% Edits still needed when they land in Session 2:
%   - "In Module 1 you ran a pre-built script" and "the Module 1 version": in the new
%     structure that was the Session 1 first-commands exercise.
%   - Step 1 "Write a CLAUDE.md for the project": participants wrote one for their own
%     project at the end of Session 1; decide whether this is a second one or reuse.
%   - Per the plan's "Session folders (decided)", this exercise runs in the Session 1
%     workspace (my_workspace/session_1), so the frame should say where to work.
%   - The old clock comments ([48:00-70:00], [70:00-75:00]) are stale.

% ══════════════════════════════════════════════════════════════════
%  EXERCISE  [48:00–70:00]
% ══════════════════════════════════════════════════════════════════

\begin{frame}{Exercise: Build the County Panel From Scratch}
  In Module 1 you ran a pre-built script. Now \textbf{you} build it.

  \vspace{0.3em}
  \textbf{Same data, fresh start.} Work in pairs using plan mode.

  \vspace{0.5em}
  \begin{enumerate}
    \item Write a \texttt{CLAUDE.md} for the project (use the template)
    \item Switch Claude Code to \textbf{plan mode}
    \item Ask Claude Code to build a county-year panel that merges election returns with IPUMS census demographics
    \item \textbf{Before approving the plan}, you will need to:
      \begin{itemize}
        \item Read the codebooks to understand variable codings\\
          {\small (What does \texttt{EDUC $\geq$ 10} mean? What does \texttt{RACE == 1} code?)}
        \item Decide how to define key variables\\
          {\small (Two-party vote share? College = bachelor's or higher?)}
        \item Specify the merge: which keys, which join type, expected row counts
      \end{itemize}
    \item Review the plan, reject or revise as needed, then approve
    \item Run the script and verify the output
  \end{enumerate}
\end{frame}

\begin{frame}{Evaluating the Plan}
  When Claude Code proposes its plan, ask yourself:

  \vspace{0.5em}
  \begin{itemize}
    \item Is it reading the \textbf{codebooks} or just guessing at variable meanings?
    \item How is it coding \texttt{EDUC}, \texttt{RACE}, \texttt{HISPAN}?\\
      {\small Do those match \textit{your} definitions?}
    \item What join type is it using? Inner, left, or outer?\\
      {\small What happens to counties that appear in one dataset but not the other?}
    \item Is it filtering \texttt{mode == "TOTAL"} for election data?
    \item Does it handle \texttt{COUNTYFIP == 0} (unidentified counties)?
  \end{itemize}

  \vspace{0.5em}
  \alert{This is where your domain knowledge matters.}\\
  Claude Code can write the code, but you decide what the variables mean.
\end{frame}

% ══════════════════════════════════════════════════════════════════
%  DEBRIEF  [70:00–75:00]
% ══════════════════════════════════════════════════════════════════

\begin{frame}{Debrief}
  \begin{itemize}
    \item What variable coding decisions did you make? Did Claude Code's default match yours?
    \item Did anyone reject or modify the plan? What did you change?
    \item How did your panel compare to the Module 1 version?
  \end{itemize}

  \vspace{0.8em}
  \begin{block}{The Lesson}
    \textbf{Writing a CLAUDE.md} forces you to be explicit about conventions you have kept implicit.

    \vspace{0.3em}
    \textbf{Reviewing plans} forces you to think about approach \textit{before} you see output.

    \vspace{0.3em}
    \textbf{Domain knowledge} is key for preventing errors.
  \end{block}
\end{frame}

